\hypertarget{cohort-and-expression}{%
\subsubsection{Cohort and expression}\label{cohort-and-expression}}

The 322 designs span 48 teams. \texttt{pb\_design\_class} (ProteinBase
modality call) labels 81 as miniprotein, 33 as nanobody, 33 as peptide,
10 as scFv, and 164 as \emph{Other}. The \emph{Other} bucket mostly
catches de novo designs the classifier could not bin. Expression rate
is 79\% (255/322; column \texttt{is\_expressed}). Expression rate varies
across method families, but denominators are too small for a defensible
test. Raw counts are in {[}Figure 2A{]}.

\hypertarget{binding}{%
\subsubsection{Binding}\label{binding}}

Nine designs cleared the binder bar (\texttt{is\_binder} =
\texttt{binding\_strength\ ∈\ \{Strong,\ Medium,\ Weak\}}). One reached
\emph{Strong} (KD \textless{} 100 nM), seven landed at \emph{Medium},
one at \emph{Weak}. The Strong design is \texttt{small-vole-maple}
(team \texttt{aryan-chandak}, method \texttt{ORBIT}, a 100-aa
miniprotein), KD = 26 nM (95\% CI from replicates: 23--28 nM, n=2 SPR
fits). The next four binders by KD are 85 nM
(\texttt{deep-bear-clay}, \texttt{kai-yi}, AF2 hallucination + ADFlip),
165 nM (\texttt{jade-cobra-cypress}, same team), 179 nM
(\texttt{young-panda-rose}, \texttt{mike-minson}, ProOS) and 186 nM
(\texttt{noble-lion-fern}, \texttt{guanjiaweitaskin}, method label not
provided). The remaining four sit in the 200--500 nM range. See
{[}Figure 2C{]} for the KD ladder.

\hypertarget{methods}{%
\subsubsection{Methods}\label{methods}}

The submitted designs span 28 named method families. The largest are
RFdiffusion (n=41), Boltz + ProteinMPNN (n=31), Mosaic (n=25), BoltzGen
(n=22), MoPPIt (n=20), LFM2 (n=14), and AF-based (n=13). BindCraft 2,
ORBIT, pXdesign, and 18 other families sit at n=7. Every design has a
matched method (\texttt{method\_family} is non-null for all 322).
Per-method hit rates are 0--14\% across the named families. With 1--3
binders per family, no per-method difference is defensible as
significant. Pacesa Lab's three binders (\texttt{amber-bat-maple},
\texttt{golden-dove-vine}, \texttt{rough-eagle-onyx}) were all
generated by BindCraft 2.

\hypertarget{modality}{%
\subsubsection{Modality}\label{modality}}

All nine binders bin as \emph{Other} in \texttt{pb\_design\_class}.
Submitter-declared modality labels eight as minibinder and one as
bispecific (\texttt{young-panda-rose}, 191 aa, ProOS). Lengths span
100--250 aa. No peptide, nanobody, or scFv design crossed the binder
bar at the default-buffer condition. The single \emph{Weak} hit is
\texttt{amber-bat-maple} (165 aa, BindCraft 2, KD = 3.2 µM).

\hypertarget{novelty}{%
\subsubsection{Novelty}\label{novelty}}

Six of nine binders sit at \textless{} 30\% sequence identity to their
closest AFDB50 neighbour (\texttt{pb\_seqidentity\_afdb50}). The three
remaining binders are antibody-derived and closer to known scaffolds
(44--62\% identity), consistent with framework-borrowing. The Strong
hit is at 21\% identity to its closest AFDB50 match.

\hypertarget{caveats}{%
\subsubsection{Caveats}\label{caveats}}

\begin{itemize}
\tightlist
\item
  9 binders. Per-method and per-modality differences are descriptive,
  not claims.
\item
  79\% expression rate refers to the 322-design ProteinBase batch.
\item
  Default-buffer screen only. A second batch in Zn²⁺-loaded buffer
  will be reported separately and may shift outcomes for designs
  targeting the holo state.
\item
  No specificity or developability data is included here; those panels
  will be reported separately.
\item
  Three teams claim binding to the disordered N-terminal tail
  (residues 12--39). No published prior art exists for binding that
  region. Treat as a novelty claim AND a specificity concern until the
  Zn²⁺ rerun and specificity panel land.
\end{itemize}
